\documentclass[11pt]{article}
\usepackage[margin=1in]{geometry}
\usepackage[T1]{fontenc}
\usepackage{lmodern,amsmath,booktabs,graphicx}
\usepackage[hidelinks]{hyperref}
\usepackage{xurl}
\setlength{\parindent}{0pt}
\setlength{\parskip}{6pt}
\setlength{\emergencystretch}{3em}
\graphicspath{{simulation_output/estimator_max_error_20261003/}}
\title{Maximum-Error Diagnosis of Pose-Initialized LiDAR Perception and NRMM Estimation}
\author{Pan Dynamics Collision Avoidance Research}
\date{October 3, 2026}
\begin{document}
\maketitle

\section{Finding and frame identification}
The largest position error is driven primarily by biased target-position
measurements. The geometric loss pulls three-dimensional body returns toward
the edges of a two-dimensional footprint, rewarding a translated and rotated
rectangle even when initialized at the true pose. The recent coverage weighting
also weakened the lateral support residual and increased that bias. The NRMM
observer then inherits the bias and develops a transient velocity error.
The separate 20 m recording has an additional, dominant startup velocity error.

The labels 15 m and 20 m refer to the initial scene gap, not the distance at
each evaluation frame. Frame indices below are zero based. Both blackout
versions have the same all-record maxima as their normal counterparts; the
maxima precede the camera interruption at 3.0--3.3 s.

\begin{center}\small
\begin{tabular}{@{}llrrrr@{}}\toprule
Recording & Maximum type & Frame & Time (s) & Sim. frame & Error (m) \\\midrule
15 m & Estimated position & 109 & 2.18 & 1374 & 0.55782 \\
15 m & Raw fitted position & 275 & 5.50 & 1540 & 0.51154 \\
20 m & Estimated position & 11 & 0.22 & 1674 & 0.35801 \\
20 m & Estimated position after 2 s & 269 & 5.38 & 1932 & 0.30411 \\
20 m & Raw fitted position & 297 & 5.94 & 1960 & 0.39778 \\
\bottomrule\end{tabular}
\end{center}

For the 15 m recording after 2 s, lateral error accounts for 97.83 percent
of raw position squared error and 95.86 percent of estimated position squared
error. The corresponding fractions at 20 m are 96.90 and 88.07 percent.
This directs the geometry investigation toward lateral center and yaw, rather
than a scalar range offset.

\section{Largest estimated-error frame: exact decomposition}
At frame 109, the actual relative center is about 4.80 m ahead. In the same
CARLA ego axes, the saved positions and their errors are
\begin{center}\small
\begin{tabular}{@{}lrrrr@{}}\toprule
 & $x$ (m) & $y$ (m) & $e_x$ (m) & $e_y$ (m) \\\midrule
True target center & 4.79886 & -0.01494 & 0 & 0 \\
Fitted observation & 4.87099 & -0.43500 & 0.07213 & -0.42005 \\
Published estimate & 4.93554 & -0.55576 & 0.13668 & -0.54082 \\
\bottomrule\end{tabular}
\end{center}
The vector identity is
\[
 \underbrace{\begin{bmatrix}0.13668\\-0.54082\end{bmatrix}}_{\text{estimated error}}
 =\underbrace{\begin{bmatrix}0.07213\\-0.42005\end{bmatrix}}_{\text{current fitted error}}
 +\underbrace{\begin{bmatrix}0.06455\\-0.12076\end{bmatrix}}_{\text{estimate minus current fit}}.
\]
These vectors are an accounting identity, not additive causal percentages.
The fitted observation has 0.42620 m error. Its heading error is
$-6.1372^\circ$, while the predicted heading error is only about
$-3.23^\circ$. With the declared 5 m length, the fitted rear-face midpoint
error and the yaw lever arm explain its lateral center error exactly:
\[
 \Delta c_y=\Delta r_{\mathrm{rear},y}
 +\frac L2\{\sin\psi_{\rm fit}-\sin\psi_{\rm true}\}
 =-0.15278-0.26727=-0.42005\ \mathrm{m}.
\]
Thus a modest false yaw creates a large center displacement. This identity
does not independently assign a physical cause to each term.

\section{Why the point-cloud objective prefers an inaccurate pose}
The frame contains 999 accepted target points; all are retained by the
two-stage fit. At the true pose, all lie inside the declared rectangular
footprint, and 721 points (72.17 percent) are more than 0.15 m from its edges.
The image and the $x$--$z$ distribution show substantial elevated body-surface
returns, consistent with trunk, window and roof surfaces. Their projection
inside a footprint is physically expected; it is not evidence that the
rectangle should move until they touch an edge.

The present nearest-face term nevertheless penalizes every projected point
according to its distance to the closest of the four footprint edges, capped
at 0.5 m. The separate containment and support-midpoint terms are too weak to
counter the reduction obtained by rotating and shifting the rectangle.
Using exactly the same 999 refinement points gives:
\begin{center}\small
\begin{tabular}{@{}lrr@{}}\toprule
Cost component (m$^2$) & True pose & Saved fitted pose \\\midrule
Capped boundary closeness & 0.116844 & 0.092531 \\
Weight-20 containment & 0 & 0.002078 \\
Coverage-weighted support midpoint & 0.000167 & 0.003624 \\
Total & 0.117011 & 0.098233 \\
\bottomrule\end{tabular}
\end{center}
The inaccurate pose improves the stated objective by about 16 percent.
Footprint-interior points contribute 0.02138 m$^2$ of the 0.02223 m$^2$
point-loss reduction before the support penalty. Increasing numerical effort
cannot repair a loss whose lower values reward this geometric bias.

\begin{figure}[p]
\centering\includegraphics[width=\textwidth]{maximum_frame_geometry.pdf}
\caption{Frame 109: recorded image, selected box, point cloud and fitted
rectangles. The green rectangle uses the true pose and declared dimensions;
it is not a separately measured true collision extent. The objective assigns
a lower cost to the inaccurate fit. Height is relative to the declared ego
origin. All plotted data come from the saved frame.}
\end{figure}

\subsection{Initial pose, local convergence and support weighting}
Several controlled changes separate the possible causes:
\begin{center}\small
\begin{tabular}{@{}lr@{}}\toprule
Frame 109 geometry control & Position error (m) \\\midrule
Current point cloud, saved predicted seed & 0.42620 \\
Same method initialized at the true pose & 0.42620 \\
Independent derivative-free minimization of the same refinement cost & 0.42479 \\
Previous objective, identical current point cloud and predicted seed & 0.20278 \\
Unit lateral support residual, independent minimization & 0.18755 \\
Lower half of points by measured height, saved predicted seed & 0.10275 \\
Points within 0.15 m of true footprint edges, diagnostic oracle & 0.09523 \\
\bottomrule\end{tabular}
\end{center}
The saved refinement reports nonconvergence after 27 iterations, but an
independent converged solve changes its position error by only 1.4 mm.
Starting at the true pose returns essentially the same biased solution.
Neither iteration count nor the extrapolated seed is the principal cause
at this frame. Other seeds can enter an opposite-yaw solution with a similar
large error; global uniqueness is not claimed.

The preceding implementation's selected lateral support residual had unit
weight. In the new method, observed lateral coverage at this fit is 0.87792,
so the residual weight is
$w=(0.87792-0.8)/0.2=0.38962$. Its squared-cost coefficient is only
$w^2=0.1518$ of the earlier unit coefficient. Correcting seed-dependent
support selection therefore also changed the objective strength. This
weakening is a concrete contributor to the recent regression, not a necessary
consequence of using the predicted pose only for initialization. Restoring
unit lateral support in a diagnostic control reduces the bias, but does not
remove the underlying body-surface/footprint mismatch.

The frame with the largest raw error, 275, has the same failure mechanism:
0.51154 m position error, $-9.0711^\circ$ yaw error, and 79.98 percent of
points farther than 0.15 m inside the true-pose footprint. A true-pose seed
still yields 0.45307 m error. The lower-height-half control gives 0.11751 m.

\subsection{Association, coordinates and timing checks}
The selected image box reaches pixel 639.36 in a 640-pixel-high image, so
the close target is clipped at the image boundary. However, a diagnostic
association using all native returns inside the true-pose declared footprint
plus 0.15 m, above ego height zero, adds 77 points and still gives 0.39445 m
error. This control uses actor truth and is not semantic instance labeling.
It shows that enlarging the image selection alone does not remove the bias.
The data provide no evidence of substitution by a parked vehicle at this frame.

Every selected frame's cached cloud agrees with the raw cloud transformed
through the fixed LiDAR-to-ego extrinsic to within $1.5\times10^{-14}$ m.
The common estimated world pose cancels from this relative-coordinate
transformation. All 300 frame identities match across native sensors,
scoring truth and prediction records in both recordings. The large error is
not explained by a frame-index shift or a world-to-ego translation mistake.
Declared dimensions remain approximate; their separate contribution to the
remaining longitudinal bias has not been isolated.

\section{How the estimator adds a transient to the measurement bias}
Fifteen actual MATLAB replays hold the saved target measurement streams
fixed, changing one input source or initialization assumption at a time.
Both baseline replays reproduce all common saved output fields exactly over
600 frames. All 4500 counterfactual publications are finite. These controls
isolate observer response; they do not rerun the changed observer state
through the perception feedback loop.

At the baseline maximum frame 109:
\begin{center}\small
\begin{tabular}{@{}lrr@{}}\toprule
Diagnostic input control & Frame 109 error (m) & Position RMSE after 2 s (m) \\\midrule
Original measurements & 0.55782 & 0.46819 \\
True target lateral position only & 0.13703 & 0.09568 \\
True target longitudinal position only & 0.54421 & 0.45998 \\
True target position, both coordinates & 0.06101 & 0.03890 \\
True ego inputs & 0.55309 & 0.46809 \\
True initial target velocity & 0.55446 & 0.46779 \\
True target position and true ego inputs & 0.03920 & 0.03877 \\
\bottomrule\end{tabular}
\end{center}
The dominant removable source is target lateral-position bias. Correcting
ego inputs or the initial target velocity alone has little effect at this
later peak. The remaining few centimetres under oracle inputs are not fully
separated into model, initialization, sampled propagation and planar-model
effects in this study.

Around frames 80--100, the fitted lateral position moves persistently away
from truth. The observer initially lags and then overshoots. Its target
velocity state acquires a false lateral component: at frame 109, target
inertial lateral velocity error is about $-0.644$ m/s. The position innovation
gain is 7.29765 s$^{-1}$, so the current 0.12076 m difference between fitted
and estimated lateral position contributes $+0.88128$ m/s to the instantaneous
position correction. Relative velocity and ego-rotation terms contribute
$-0.64469$ and $+0.00287$ m/s. This predicts recovery after the peak; these
instantaneous terms are not an exact integrated decomposition of the past.

The runtime publishes the current predicted state before integrating the
current observation over the next sample interval. Replacing only frame
109's lateral observation with truth therefore leaves frame 109's published
error unchanged at 0.55782 m. Replacing only frame 108's observation reduces
it to 0.49667 m. The extra 0.12076 m lateral displacement relative to the
current fit reflects prior history and observer dynamics, rather than an
additional error introduced by the frame 109 optimizer.

\begin{figure}[p]
\centering\includegraphics[width=\textwidth]{error_sources_trace.pdf}
\caption{Complete recorded traces. At 15 m the lateral fit bias persists
even after the predicted heading returns near truth. At 20 m, changes
between geometric fitting solutions produce lateral oscillations and delayed
observer extrema. True ego-input replays nearly overlap the original result;
true target-position inputs remove most later error. All oracle curves are
diagnostic, not deployable measurements.}
\end{figure}

\section{The 20 m recording has two different peak mechanisms}
Its all-record maximum occurs at frame 11, 0.22 s: the position error vector
is $(0.31814,-0.16418)$ m. The current fitted observation has only 0.01748 m
longitudinal error but 0.16348 m lateral error. The observer's initial
co-moving target velocity is approximately 8.0687 m/s forward, while actual
target forward velocity is 4.7674 m/s. The full initial velocity error is
3.3016 m/s. This creates a longitudinal startup transient that is distinct
from the later lateral fitting problem.

\begin{center}\small
\begin{tabular}{@{}lr@{}}\toprule
20 m diagnostic control & Error at original frame 11 (m) \\\midrule
Original measurements and initialization & 0.35801 \\
True target position inputs & 0.29427 \\
True initial target velocity & 0.16052 \\
True target positions and true initial target velocity & 0.02072 \\
True ego inputs & 0.36001 \\
\bottomrule\end{tabular}
\end{center}
Later, the estimated maximum after 2 s is frame 269, 5.38 s. Its 0.30411 m
estimate error occurs when the current raw fit error is only 0.06972 m,
following earlier positive lateral fits. This is another history-dependent
observer extremum. The largest raw fit error instead occurs at frame 297:
0.39778 m and $-5.5959^\circ$ yaw error. Initializing that frame at truth
still gives 0.39778 m. The lower-height-half control reduces its raw error
to 0.06120 m. After 2 s over the entire recording, true position inputs
reduce estimated position RMSE from 0.13940 to 0.01729 m.

\section{Supported next change and limits of the current evidence}
The first priority is a visibility-aware geometric objective: retain
containment information from all valid target returns, but apply proximity
residuals to points supporting observed footprint faces or silhouette edges.
Three-dimensional surface orientation or height structure can help separate
elevated body returns from footprint-boundary support. All three pose
variables should remain free, with extrapolation supplying only the initial
value. The lateral support term also needs deliberate strength and visibility
semantics, rather than weakening merely because candidate orientation changes.

As a whole-record diagnostic, selecting the lower half of points by measured
height while holding every original predicted seed fixed produces 300 valid
fits per recording. Raw position RMSE after 2 s changes from 0.45718 to
0.11143 m at 15 m and from 0.20377 to 0.03460 m at 20 m. This independently
supports the point-population diagnosis beyond the single worst frame.
It is not a validated deployment rule: the split retains a longitudinal
bias, and arbitrary viewing directions, occlusions and body shapes have not
been evaluated. These are fixed-seed point ablations, not new closed-loop
perception/estimator results.

The second priority is target velocity initialization from time-aligned
early observations and ego motion, rather than assuming the target initially
shares ego velocity. Observer gain changes or more geometry iterations are
lower priorities: the current observer cannot identify a persistent lateral
measurement offset solely from that same biased measurement stream, and
independent geometric optimization preserves the bias.

This study changes no production algorithm or gain. It uses the same two
previously seen straight-route recordings, native LiDAR, frozen actual
detector selections and the original native ego preprocessing, including
its ideal-compass limitation. No CARLA connection or detector inference
was performed. Truth is deliberately used in labeled diagnostic controls;
those controls are not claims of measured-input performance. No global
optimality, generalization or new observer error certificate is established.

\section{Reproduction and artifacts}
The analyzed perception version is
\path{4851f4bdc3f7047cec7ad8664d3c98e8a753d0e5}; the actual MATLAB runtime is
the isolated existing snapshot
\path{26cd6f3998d17fd49f1cd9254245805d17f4d3ae}. The matched previous-objective
comparison uses only the geometric function from
\path{447056c8f2092e0691a27da15525436342e409b2}, without restoring any old
production path. Original LiDAR/ego-sensor seed pairs remain
20261002/20261003 and 20261003/20261004; the interval is 0.02 s.

Five saved point fits were reproduced to $10^{-10}$ tolerance, with 50
selected-frame geometry controls, five matched previous-objective controls,
600 whole-record point-selection fits and 15 MATLAB replays. The independent
derivative-free optimizer passed a known-quadratic minimum check. Existing
production unit suites were not rerun for this report-only change.

Compact measurements, provenance and validation are in
\path{report/ESTIMATOR_MAX_ERROR_DIAGNOSIS_20261003/}. Local diagnostic
scripts, input substitutions, raw traces, figure sources and the derived PDF
are in \path{simulation_output/estimator_max_error_20261003/}. Recorded
datasets, runtime copies, historical comparison code, dependencies and
generated figures/binaries remain outside Git. Identified archive exports
preserve the supporting study artifacts and original evidence references.
\end{document}
